\section{Relational Trust}
\label{p13:sec:relational}

\noindent
The mechanism of \xref{p13:sec:mechanism} described how trust is generated between a promisor and a receiver who are present to one another: the structural form, the costly self-binding, the wagered leap, all presuppose two parties in some direct relation. But a great part of the trust on which human dealing actually runs is not of this direct kind. We trust people we have never met, on the strength of others we do trust; trust travels along the links of a relational network, reaching where no direct relation reaches. This section is about that transmitted trust\,---\,\emph{relational trust}\,---\,its mechanism, its peculiar liability to a failure the direct case does not suffer, and the formal structure it shares, the section will argue, with the dynamics of an open quantum system. It is the chapter in which the problem of winning a family's trust, posed in the Prelude, finds the resource that bears most directly upon it; and it is the chapter in which the crisis of the symbol receives a new and precise specification.

\subsection{The mechanism of transmission: the recommendation and the chain}
\label{p13:sub:rel-chain}

The clearest instance of relational trust is the letter of recommendation, and it repays close mechanical attention, for it shows what is really at work and what is not. The naive view is that a recommendation works by the \emph{content} of what the recommender says\,---\,by the predicates he applies to the candidate, the praise he heaps, the qualities he affirms. But this is very largely not what carries the trust. The adjectives carry almost no information; everyone is described in superlatives, and a committee that has read a thousand letters extracts nothing from the thousand-and-first profession of excellence (\xref{p13:sub:sympt-universal}). What actually works is not the content but the \emph{structure}: a committee $C$ trusts the judgment of a recommender $A$, with whom it has a relation of trust; the recommender $A$ has a real working coupling with the candidate $B$, knows $B$ from the inside, has the embodied knowledge of what $B$ can do; and trust is transmitted along the chain $C\!-\!A\!-\!B$, $C$ coming to trust $B$ \emph{through} its trust in $A$. The recommendation works as a conductor, not as a description: it carries $C$'s existing trust in $A$ across the link $A\!-\!B$ to reach $B$, whom $C$ has no other way to reach.

What this reveals is that trust, in the relational case, is not a property of any single party at all. It is not something in $B$, nor in $A$, nor in $C$; it is a property of the \emph{structure of the relation} among them\,---\,a topological property of the network, residing in the pattern of couplings rather than in any node. $C$ has no direct evidence about $B$ (the evidentialist channel cannot reach), and no direct contact with $B$ (the deontological reading of character cannot reach); yet $C$ can trust $B$, and the trust exists nowhere but in the chain. Relational trust is thus a third mechanism, irreducible to the evidential and the characterological alike, by which trust reaches across a gap that neither evidence nor the direct reading of character could span. And it bears directly on the predicament of the Prelude: the trust a family extends to a stranger who would marry into it runs, very largely, not through anything the stranger can show them directly, but through the chain that connects him to them by way of the one they already love and trust\,---\,the one who loves him is, in the structure of the network, his recommender, and the trust reaches him along that link or not at all.

\subsection{Relational trust spans the three registers}
\label{p13:sub:rel-threeregisters}

Relational trust is not lodged in one register but spans all three, and this is the key to both its power and its specific fragility. There is an imaginary component: the image of $B$ that travels along the chain, the idealised or distorted picture $A$'s account installs in $C$. There is a symbolic component: the letter, the word, the chain itself as a transmissible structure of signs. And there is a real component, which is the heart of the matter: the real working coupling between $A$ and $B$, $A$'s embodied, unfakeable knowledge of what $B$ is and can do, the thing the whole apparatus exists to transmit. Genuine relational trust is the case in which the real component leads\,---\,in which what travels the chain is a real coupling, vouched along a structure of signs, forming in $C$ an expectation grounded in something that actually exists. The chain is the symbolic conductor; what it ought to conduct is the real.

\subsection{The dislocation of vehicle from substance}
\label{p13:sub:rel-dislocation}

Here is the new specification of the crisis of the symbol that the section promised, and it follows from the fact just established\,---\,that relational trust spans the registers, and that the registers can be \emph{prised apart}. The substance relational trust ought to transmit is the real component: $A$'s real coupling with $B$, the unfakeable knowledge that is the one thing of value in the whole transaction. But the vehicle by which it is transmitted is symbolic: a letter, a word, a writable and circulable and archivable sign. And the alienation that destroys relational trust occurs in the dislocation of the vehicle from the substance\,---\,the symbolic form preserved intact while the real substance is hollowed out of it.

A forged recommendation is exactly this dislocation. In the symbolic register it is impeccable: correctly formatted, well phrased, duly signed, transmitted along the chain $C\!-\!A\!-\!B$ in perfect order. Viewed from the symbolic register alone it is indistinguishable from a genuine one. But its real component\,---\,$A$'s actual knowledge of $B$\,---\,is empty: $A$ wrote it not because he embodily knows $B$ to be good but because he owed a favour, or moves in the same circle, or could not decently refuse. The sign is transmitted, and what the sign was supposed to carry is not present in it. This is the incarnation, in relational trust, of the central pathology the series has named the failure of settlement, the usurpation by the symbol: a sign that announces a coupling already vouched-for when no real coupling underwrites it, a letter that cashes a knowledge that was never had. The forged recommendation is the isomorph of the false poem and the alienated bride-price\,---\,a symbol that claims to anchor the real while the real has been drained out from beneath it.

The diagnosis of the breakdown is therefore precise, and it is the section's new contribution to the diagnosis of \xref{p13:sec:diagnosis}. Relational trust breaks down not because the mechanism is defective but because it spans the registers and the registers can be split: the symbolic transmissibility (the letter, the chain) and the real unfakeability (the actual coupling), which ought to be two faces of one thing, can be prised into two, the first retained and the second hollowed out. Forged relational trust is the product of that fissure.

\subsection{The bidirectional alienation}
\label{p13:sub:rel-bidirectional}

The word \emph{alienation} is exact here, for the damage runs in both directions along the chain, and the section must trace both, for together they explain why the whole channel closes rather than merely admitting some bad letters alongside the good.

On the side of the one vouched-for, the proliferation of forged recommendations devalues the genuine ones. The student who really does have a deep working coupling with a mentor who really knows her finds that her letter, in the symbolic register, looks exactly like the letter written from obligation; and so her real coupling can no longer be transmitted, because the channel that should carry it has been saturated with signs that carry nothing. The real thing loses its capacity to be recognised, not because it is disbelieved in particular but because the forgeries have made the whole channel unreadable. This is the precise mechanism of the crisis of the symbol, seen once more: it is not that the false drives out the true by replacing it, but that the false makes the true \emph{unrecognisable}, so that the channel as a whole is abandoned and the institution retreats to the merely countable\,---\,citation counts, rankings\,---\,signals that are barren but at least unforgeable, preferring a channel known to be poor to one that might be rich but has been poisoned.

On the side of the one who vouches, the damage is subtler and is an alienation in the fullest sense. The recommender who grows accustomed to writing letters he does not believe is alienated from his own act of vouching: his signature ceases to mean that he really knows and comes to mean only that a social obligation has been discharged. And in doing so he draws down, a little at a time, the relational capital that made relational trust possible\,---\,$C$'s trust in his judgment\,---\,spending it to discharge favours until the trust itself is exhausted. He siphons off, to meet the demands of his circle, the very thing that let his word carry trust, and when it is gone the chain through him conducts nothing. The bidirectional alienation\,---\,the genuine made unrecognisable on the one side, the capacity to vouch drawn down to exhaustion on the other\,---\,is why the failure of relational trust is not a local matter of some bad letters but the closing of a channel.

\subsection{The quantum-dynamical structure of relational trust}
\label{p13:sub:rel-quantum}

The structure that has emerged\,---\,a trust that is a property of a network rather than of its nodes, that spans registers which can be split, that fails by a dislocation of a transmissible form from an unfakeable substance\,---\,has, the section now argues, a precise formal counterpart in the dynamics of open quantum systems. The argument is made under a discipline that must be stated before it is made, and stated as firmly as the discipline upon the use of predictive coding in \xref{p13:sec:mechanism}. \emc{What is claimed is a structural isomorphism, not a physical reduction.} The paper does not assert that trust is, literally, a quantum phenomenon; to say so would be the worst kind of pseudo-scientific mysticism, and would forfeit the seriousness of everything else the paper has done. What is asserted is that the mathematical structure of relational trust and the mathematical structure of the entanglement-and-decoherence dynamics of an open quantum system are the same structure, and that the formal language of the latter can therefore describe the former with a precision no other available language affords. This is the same discipline the series has observed throughout: the quantum-dynamical account is one more level of description, a formal language for a structure, and not the truth beneath the others, nor a meta-framework that absorbs them.

Under that discipline, three correspondences may be drawn, in ascending order of depth.

The first is \emph{entanglement} as the non-local structure of relational trust. The defining feature of an entangled state is that the state of the whole cannot be decomposed into a product of the states of its parts; the information resides in the correlations among the parts, not in any part taken alone, so that a measurement of any single component does not recover it. This is exactly the structure established in \xref{p13:sub:rel-chain}: relational trust is not a property of $C$, $A$, or $B$ severally, but of the correlation $C\!-\!A\!-\!B$ itself, recoverable from no node in isolation. The non-locality of trust\,---\,that it lives in the chain and not in the links\,---\,is given, by the language of entanglement, a precise formal expression it otherwise lacks.

The second is \emph{decoherence} as the dislocation of vehicle from substance. A coherent quantum state, coupled to an environment, undergoes decoherence: the correlations that constitute its coherence are not destroyed outright, but the phase relations that made them meaningful are lost to the environment, leaving a state whose correlational shell persists while the coherence that gave it its character has drained away. This is, formally, the dislocation of \xref{p13:sub:rel-dislocation}: the symbolic transmissibility (the correlational shell, the chain of signs) is retained, while the real coupling (the phase coherence that made the correlation carry something) is lost. A forged recommendation is a \emph{decohered} trust-state: it preserves, in the symbolic register, the entire outward form of the entanglement, while the coherence it was meant to carry is gone\,---\,and, just as a decohered mixed state can be indistinguishable from a pure state under measurements that are not phase-sensitive, the forged letter is indistinguishable from the genuine under any reading that is not sensitive to the real coupling beneath. And the crisis of the symbol is, on this reading, the \emph{systemic decoherence} of relational trust at the scale of a society: an environment saturated with forged signals continually coupling to the trust-states that traverse it, draining their coherence, leaving correlational shells that no longer carry what they were meant to carry. This dovetails exactly with the predictive-coding diagnosis of \xref{p13:sub:mech-crisis}: the systemic decoherence of the trust-state and the rational collapse of the population's precision-weight on the symbolic channel are the same process described at two levels.

Here the word \emph{dynamical}, rather than merely \emph{quantum}, earns its place, for what relational trust most needs from the analogy is not a static picture of a state but an account of how the state \emph{evolves}. The dynamics of an \emph{open} quantum system\,---\,a system in continual exchange with an environment, its evolution governed by a master equation that specifies how coherence decays, and under what conditions it may be protected or even restored\,---\,is the formal language for how a trust-state is maintained, corrupted, or rebuilt over time, which a static account of entanglement could not supply. And this is the point at which the present section welds to \xref{p13:sec:holonomy}, for the geometric phase deployed there is, in its origin, a concept of exactly this quantum dynamics: the Berry phase is the phase accumulated by a quantum state under adiabatic transport. The holonomy language of \xref{p13:sec:holonomy} was therefore never a free-floating borrowing; it was, all along, drawn from the very dynamics relational trust is here shown to share. Adiabatic evolution\,---\,slow, exchanging no energy with the environment, coherence and geometric phase preserved\,---\,is the genuine, endogenously sustained trust; non-adiabatic evolution\,---\,coupled to the environment, decohering, phase lost\,---\,is the forged. The genuine-versus-forged distinction of \xref{p13:sec:justice}, the adiabatic-versus-driven distinction of \xref{p13:sec:holonomy}, and the coherent-versus-decohered distinction drawn here are one distinction in three vocabularies.

The third and most precise correspondence concerns transmission itself, and it requires a refinement that shows the analogy has been taken seriously rather than reached for loosely. Bare entanglement will not, by itself, model the \emph{transmission} of trust, for entanglement is subject to the no-communication theorem: it cannot, alone, be used to send information. But relational trust is precisely about the transmission of trust along the chain $C\!-\!A\!-\!B$. The formal counterpart of this transmission is therefore not bare entanglement but \emph{entanglement swapping}: the protocol by which, through an intermediary node $A$ already entangled with each, two parties $C$ and $B$ who share no prior entanglement are brought into an entangled relation. Entanglement swapping is, almost exactly, the quantum correspondent of the recommendation mechanism $C\!-\!A\!-\!B$: the intermediary $A$, coupled to both, is the means by which a correlation is established between $C$ and $B$ who had none. To name the correspondent as swapping rather than as bare entanglement is to mark that the analogy is structural and considered\,---\,that it tracks the specific feature of relational trust, its transmissibility through an intermediary, and not merely its non-locality.

\subsection{The standing of the analogy}
\label{p13:sub:rel-standing}

The discipline stated at the opening of \xref{p13:sub:rel-quantum} must be restated at its close, for it is the condition on which the whole construction is admissible. The quantum-dynamical language is a formal description of a shared structure, at a level of description that does not displace the others the paper has given\,---\,the mechanical account of \xref{p13:sec:mechanism}, the ethical account of \xref{p13:sec:justice}, the phenomenological account of the chain as a relation among persons. It is not a claim that trust is quantum-mechanical in its physical substrate; it is a claim that relational trust and open-system quantum dynamics instantiate one mathematical structure, and that the structure is illuminated by being seen in the language that has developed the richest tools for it. It is, in the polyphonic terms the series requires, one more irreducible epistemic position on a single object, a further voice and not the final word\,---\,and it is offered as such, a formal analogy that earns its place by the precision it brings to the non-locality, the decoherence, and the transmission of a trust that lives in the chain and not in its links.
